%%%%%%%%%%  scan idx 432  |  folio 419  %%%%%%%%%%
régulier, si, localement sur $X$, il existe un $A$-Module libre de type fini $E$, et un homomorphisme régulier surjectif : $E\longrightarrow J$.

Soit $i:Y\longrightarrow X$ une immersion de préschémas, et soit $U$ un ouvert de $X$ tel que $i(Y)\subset U$ et que $i$ soit une immersion fermée de $Y$ dans $U$. On dit que $i$ est régulière si l'idéal de $\cO_U$ définissant $i(Y)$ est régulier (condition qui ne dépend pas de $U$). Si $y\in Y$, on dit que $i$ est régulière en $y$ s'il existe un voisinage $V$ de $i(y)$ tel que l'immersion $i^{-1}(V)\longrightarrow V$ soit régulière. Pour qu'une immersion soit régulière, il est donc nécessaire et suffisant qu'elle le soit au voisinage de chaque point.

\emph{Cette définition est en conflit avec} EGA IV 16.9.2, \emph{et remplace cette dernière}. Toutefois, dans le cas noethérien, les deux définitions sont équivalentes (EGA IV 19.5.1). Une immersion qui peut localement être définie par une suite régulière est régulière; une immersion régulière est quasi-régulière (1.3).

\textbf{1.4.1.} Soient $X$ un schéma, $J$ un Idéal de type fini de $\cO_X$, $x$ un point de $\Supp(\cO_X/J)$. Nous dirons qu'une famille de sections $f_1,\ldots,f_d$ de $J$ sur un voisinage de $x$ forme un \emph{système minimum de générateurs de $J$ au voisinage de} $x$ si les $(f_i)_x$ sont une famille de générateurs de $J_x$ de nombre minimum, i.e. leurs images dans $J_x\otimes k(x)$ forment une base de cet espace vectoriel sur $k(x)$. On notera que les $f_i$ engendrent $J$ sur un voisinage de $x$.

\textbf{1.4.2.} Dire qu'une immersion est régulière, c'est encore dire qu'elle est définie dans un certain ouvert de $X$ par un idéal de type fini $J$, et que pour tout point $x\in\Supp(\cO_X/J)$, et tout système minimum de générateurs de $J$ au voisinage de $x$, le complexe de Koszul associé à cette suite est acyclique au voisinage de $x$: soit en effet $f_1,\ldots,f_d$ un système minimum de
